\subsection{Digital Filters}
\label{sec:digital-filters}

A digital filter computes each output sample $y[n]$ from a combination of past and present input and output values.
Two fundamental classes are distinguished by whether feedback is present.

\subsubsection{FIR and IIR Filters}

A \textit{Finite Impulse Response} (FIR) filter computes $y[n]$ as a weighted sum of the $M$ most recent input samples only:

\[
y[n] = \sum_{k=0}^{M} b_k \, x[n-k]
\]

Because there is no feedback, FIR filters are unconditionally stable and can achieve a \textit{linear phase} response, which preserves the waveshape of the filtered signal.
Their main disadvantage is computational cost: achieving a steep roll-off requires a large number of coefficients \cite[p.~293]{oppenheimDiscretetimeSignalProcessing2014}.

An \textit{Infinite Impulse Response} (IIR) filter adds feedback terms --- past output values --- to the computation:

\[
y[n] = \sum_{k=0}^{M} b_k \, x[n-k] \;-\; \sum_{k=1}^{N} a_k \, y[n-k]
\]

The recursive structure allows sharp frequency responses with far fewer coefficients than an equivalent FIR filter, making IIR filters well suited to resource-constrained embedded systems.
The trade-off is that stability is not guaranteed: if the poles of the transfer function lie outside the unit circle in the $z$-plane, the filter becomes unstable \cite[p.~371]{oppenheimDiscretetimeSignalProcessing2014}.
